% questions.tex —— 2014 年·广东卷·文科 · 题目内容片段（无答案）
% 由 tools/generate_projects.py 依来源总账自 2014/2014_数学真题_广东文_试题.md 生成，勿手改（改源 MD 后重跑）


\begin{gknotice}
\textbf{来源说明}：本题面依据固定源 PDF 的两页 A3 题面逐题转录，原卷印刷页码为 879—880。第 879 面为选择题 1—10、填空题 11—15 与解答题 16—18；第 880 面为解答题 19—21。全卷三处节头依次为「一、选择题」「二、填空题」「三、解答题」。

\textbf{解析归属}：固定源不附答案或解析；配套答案与解析均由本库按本题面独立复算。

\textbf{图形说明}：第 15 题的平行四边形与第 18 题折叠前后两图由项目内 TikZ 矢量源重绘；第 17 题为年龄数据表，原页无图；其余各题原页无图。
\end{gknotice}

\gkpart{一、选择题}

\begin{question}
\gkti{1}已知集合 $M=\{2,3,4\}$，$N=\{0,2,3,5\}$，则 $M\cap N=$\mbox{（　　）}

（A）$\{0,2\}$；（B）$\{2,3\}$；（C）$\{3,4\}$；（D）$\{3,5\}$。

\end{question}
\begin{question}
\gkti{2}已知复数 $z$ 满足 $(3-4\mathrm i)z=25$，则 $z=$\mbox{（　　）}

（A）$-3-4\mathrm i$；（B）$-3+4\mathrm i$；（C）$3-4\mathrm i$；（D）$3+4\mathrm i$。

\end{question}
\begin{question}
\gkti{3}已知向量 $\boldsymbol a=(1,2)$，$\boldsymbol b=(3,1)$，则 $\boldsymbol b-\boldsymbol a=$\mbox{（　　）}

（A）$(-2,1)$；（B）$(2,-1)$；（C）$(2,0)$；（D）$(4,3)$。

\end{question}
\begin{question}
\gkti{4}若变量 $x$，$y$ 满足约束条件 $\begin{cases}x+2y\leqslant 8,\\0\leqslant x\leqslant 4,\\0\leqslant y\leqslant 3,\end{cases}$ 则 $z=2x+y$ 的最大值等于\mbox{（　　）}

（A）$7$；（B）$8$；（C）$10$；（D）$11$。

\end{question}
\begin{question}
\gkti{5}下列函数为奇函数的是\mbox{（　　）}

（A）$2^x-\dfrac1{2^x}$；（B）$x^3\sin x$；（C）$2\cos x+1$；（D）$x^2+2^x$。

\end{question}
\begin{question}
\gkti{6}为了解 $1000$ 名学生的学习情况，采用系统抽样的方法，从中抽取容量为 $40$ 的样本，则分段的间隔为\mbox{（　　）}

（A）$50$；（B）$40$；（C）$25$；（D）$20$。

\end{question}
\begin{question}
\gkti{7}在 $\triangle ABC$ 中，角 $A$，$B$，$C$ 所对应的边分别为 $a$，$b$，$c$，则“$a\leqslant b$”是“$\sin A\leqslant\sin B$”的\mbox{（　　）}

（A）充分必要条件；（B）充分非必要条件；（C）必要非充分条件；（D）非充分非必要条件。

\end{question}
\begin{question}
\gkti{8}若实数 $k$ 满足 $0<k<5$，则曲线 $\dfrac{x^2}{16}-\dfrac{y^2}{5-k}=1$ 与曲线 $\dfrac{x^2}{16-k}-\dfrac{y^2}{5}=1$ 的\mbox{（　　）}

（A）实半轴长相等；（B）虚半轴长相等；（C）离心率相等；（D）焦距相等。

\end{question}
\begin{question}
\gkti{9}若空间中四条两两不同的直线 $l_1$，$l_2$，$l_3$，$l_4$，满足 $l_1\perp l_2$，$l_2\parallel l_3$，$l_3\perp l_4$，则下列结论一定正确的是\mbox{（　　）}

（A）$l_1\perp l_4$；（B）$l_1\parallel l_4$；（C）$l_1$，$l_4$ 既不垂直也不平行；（D）$l_1$，$l_4$ 的位置关系不确定。

\end{question}
\begin{question}
\gkti{10}对任意复数 $w_1$，$w_2$，定义 $w_1*w_2=w_1\overline{w_2}$，其中 $\overline{w_2}$ 是 $w_2$ 的共轭复数，对任意复数 $z_1$，$z_2$，$z_3$ 有如下四个命题：

① $(z_1+z_2)*z_3=(z_1*z_3)+(z_2*z_3)$；② $z_1*(z_2+z_3)=(z_1*z_2)+(z_1*z_3)$；

③ $(z_1*z_2)*z_3=z_1*(z_2*z_3)$；④ $z_1*z_2=z_2*z_1$。

则真命题的个数是\mbox{（　　）}

（A）$1$；（B）$2$；（C）$3$；（D）$4$。

\end{question}
\gkpart{二、填空题}

\begin{question}
\gkti{11}曲线 $y=-5\mathrm e^x+3$ 在点 $(0,-2)$ 处的切线方程为\gkblank{}。

\end{question}
\begin{question}
\gkti{12}从字母 $a$，$b$，$c$，$d$，$e$ 中任取两个不同字母，则取到字母 $a$ 的概率为\gkblank{}。

\end{question}
\begin{question}
\gkti{13}等比数列 $\{a_n\}$ 的各项均为正数，且 $a_1a_5=4$，则 $\log_2a_1+\log_2a_2+\log_2a_3+\log_2a_4+\log_2a_5=$\gkblank{}。

\end{question}
\begin{question}
\gkti{14}在极坐标系中，曲线 $C_1$ 和 $C_2$ 的方程分别为 $2\rho\cos^2\theta=\sin\theta$ 和 $\rho\cos\theta=1$，以极点为平面直角坐标系的原点，极轴为 $x$ 轴的正半轴，建立平面直角坐标系，则曲线 $C_1$ 和 $C_2$ 的交点的直角坐标为\gkblank{}。

\Needspace{0.4\textheight}
\end{question}
\begin{question}
\gkti{15}如图，在平行四边形 $ABCD$ 中，点 $E$ 在 $AB$ 上，且 $EB=2AE$，$AC$ 与 $DE$ 交于点 $F$，则 $\dfrac{\triangle CDF\ \text{\text{的周长}}}{\triangle AEF\ \text{\text{的周长}}}=$\gkblank{}。
\nopagebreak[4]
\begin{center}% ============================================================
% f2014_guangdong_arts_q15_parallelogram.tex
% 2014 广东卷（文科）第 15 题：平行四边形 ABCD 中 E 在 AB 上，
% EB=2AE，DE 与 AC 交于 F，求 △CDF 与 △AEF 的周长比。
% 来源 PDF：普通高考/2014/2014广东文.pdf
% 原页：第 1 页（A3 横排）第 2 栏，印面 879，第 15 题题干下方
% 说明：AE:EB=1:2，故 AE=AB/3；图中 A、E、B 依次在底边上，
% D 在 A 上方偏右，C 在 B 上方偏右，DE 近于竖直，AC 为对角线。
% ============================================================
\begin{tikzpicture}[>=Stealth, line width=0.7pt, font=\small,
  x=1cm, y=1cm, line cap=round, line join=round]

  \coordinate (A) at (0,0);
  \coordinate (E) at (2,0);
  \coordinate (B) at (6,0);
  \coordinate (D) at (1.6,3);
  \coordinate (C) at (7.6,3);
  \coordinate (F) at (1.9,0.75);

  \draw (A) -- (B) -- (C) -- (D) -- cycle;
  \draw (D) -- (E);
  \draw (A) -- (C);

  \fill (A) circle (1.1pt) node[below left] {$A$};
  \fill (E) circle (1.1pt) node[below] {$E$};
  \fill (B) circle (1.1pt) node[below right] {$B$};
  \fill (D) circle (1.1pt) node[above left] {$D$};
  \fill (C) circle (1.1pt) node[above right] {$C$};
  \fill (F) circle (1.1pt) node[right] {$F$};
\end{tikzpicture}
\end{center}
\nopagebreak[4]
\end{question}
\gkpart{三、解答题}

\begin{question}
\gkti{16}已知函数 $f(x)=A\sin\left(x+\dfrac\pi3\right)$，$x\in\mathbf R$，且 $f\left(\dfrac{5\pi}{12}\right)=\dfrac{3\sqrt2}{2}$。

（1） 求 $A$ 的值；

（2） 若 $f(\theta)-f(-\theta)=\sqrt3$，$\theta\in\left(0,\dfrac\pi2\right)$，求 $f\left(\dfrac\pi6-\theta\right)$。

\end{question}
\begin{question}
\gkti{17}某车间 $20$ 名工人年龄数据如下表：

\begin{center}
\begin{tabular}{ll}
\toprule
年龄（岁） & 工人数（人） \\
\midrule
$19$ & $1$ \\
$28$ & $3$ \\
$29$ & $3$ \\
$30$ & $5$ \\
$31$ & $4$ \\
$32$ & $3$ \\
$40$ & $1$ \\
合计 & $20$ \\
\bottomrule
\end{tabular}
\end{center}

（1） 求这 $20$ 名工人年龄的众数与极差；

（2） 以十位数为茎，个位数为叶，作出这 $20$ 名工人年龄的茎叶图；

（3） 求这 $20$ 名工人年龄的方差。

\Needspace{0.4\textheight}
\end{question}
\begin{question}
\gkti{18}如图 1，四边形 $ABCD$ 为矩形，$PD\perp$ 平面 $ABCD$，$AB=1$，$BC=PC=2$，作如图 2 折叠，折痕 $EF\parallel DC$。其中点 $E$，$F$ 分别在线段 $PD$，$PC$ 上，沿 $EF$ 折叠后点 $P$ 叠在线段 $AD$ 上的点记为 $M$，并且 $MF\perp CF$。
\nopagebreak[4]
\begin{center}% ============================================================
% f2014_guangdong_arts_q18_fold.tex
% 2014 广东卷（文科）第 18 题：矩形 ABCD 中 PD⊥平面 ABCD，
% AB=1，BC=PC=2，沿 EF（EF∥DC，E∈PD，F∈PC）折叠，P 落在
% 线段 AD 上的 M 处，且 MF⊥CF。左为折叠前（图 1），右为折叠后（图 2）。
% 来源 PDF：普通高考/2014/2014广东文.pdf
% 原页：第 1 页（A3 横排）第 3 栏下部，印面 879，第 18 题题干下方
% 说明：图 2 依原页保留折叠前 P 点及 P—E、P—F 两段（示意折痕两侧），
% 并以虚线表示被遮挡的棱 MD、DC、DE、EC；三棱锥 M-CDE 的 ME、MC 为实线。
% ============================================================
\begin{tikzpicture}[>=Stealth, line width=0.7pt, font=\small,
  x=1cm, y=1cm, line cap=round, line join=round]

  % ---------------- 图 1：折叠前 ----------------
  \coordinate (A1) at (0,4.6);
  \coordinate (B1) at (3.4,4.6);
  \coordinate (C1) at (3.72,1.0);
  \coordinate (D1) at (0.3,1.0);
  \coordinate (P1) at (-1.0,-1.0);

  \draw (A1) -- (B1) -- (C1) -- (D1) -- cycle;
  \draw (P1) -- (D1);
  \draw (P1) -- (C1);

  \fill (A1) circle (1.1pt) node[above left] {$A$};
  \fill (B1) circle (1.1pt) node[above right] {$B$};
  \fill (C1) circle (1.1pt) node[right] {$C$};
  \fill (D1) circle (1.1pt) node[above left] {$D$};
  \fill (P1) circle (1.1pt) node[below left] {$P$};
  \node at (1.7,-1.35) {图 1};

  % ---------------- 图 2：折叠后 ----------------
  \coordinate (A2) at (6.2,4.6);
  \coordinate (B2) at (9.6,4.6);
  \coordinate (C2) at (9.92,1.0);
  \coordinate (D2) at (6.5,1.0);
  \coordinate (P2) at (5.2,-1.0);
  \coordinate (M2) at (6.3164,3.2032);
  \coordinate (E2) at (6.175,0.5);
  \coordinate (F2) at (8.74,0.5);

  % 矩形边框（DC 被折起部分遮挡）
  \draw (A2) -- (B2) -- (C2);
  \draw (A2) -- (M2);
  \draw[dashed] (M2) -- (D2) -- (C2);

  % 被遮挡的棱
  \draw[dashed] (E2) -- (D2);
  \draw[dashed] (E2) -- (C2);

  % 折痕与折叠后的三角形 MEF
  \draw (E2) -- (F2);
  \draw (F2) -- (C2);
  \draw (M2) -- (E2);
  \draw (M2) -- (F2);
  \draw (M2) -- (C2);

  % 折叠前 P 的位置（示意）
  \draw[line width=0.55pt] (P2) -- (E2);
  \draw[line width=0.55pt] (P2) -- (F2);

  \fill (A2) circle (1.1pt) node[above left] {$A$};
  \fill (B2) circle (1.1pt) node[above right] {$B$};
  \fill (C2) circle (1.1pt) node[right] {$C$};
  \fill (D2) circle (1.1pt) node[above right, inner sep=2pt] {$D$};
  \fill (M2) circle (1.1pt) node[left] {$M$};
  \fill (E2) circle (1.1pt) node[left] {$E$};
  \fill (F2) circle (1.1pt) node[below right] {$F$};
  \fill (P2) circle (1.1pt) node[below left] {$P$};
  \node at (7.9,-1.35) {图 2};
\end{tikzpicture}
\end{center}
\nopagebreak[4]
（1） 证明：$CF\perp$ 平面 $MDF$；

（2） 求三棱锥 $M-CDE$ 的体积。

\end{question}
\begin{question}
\gkti{19}设各项均为正数的数列 $\{a_n\}$ 的前 $n$ 项和为 $S_n$，且 $S_n$ 满足 $S_n^2-(n^2+n-3)S_n-3(n^2+n)=0$，$n\in\mathbf N^*$。

（1） 求 $a_1$ 的值；

（2） 求数列 $\{a_n\}$ 的通项公式；

（3） 证明：对一切正整数 $n$，有 $\dfrac1{a_1(a_1+1)}+\dfrac1{a_2(a_2+1)}+\cdots+\dfrac1{a_n(a_n+1)}<\dfrac13$。

\end{question}
\begin{question}
\gkti{20}已知椭圆 $C:\dfrac{x^2}{a^2}+\dfrac{y^2}{b^2}=1$（$a>b>0$）的一个焦点为 $(\sqrt5,0)$，离心率为 $\dfrac{\sqrt5}{3}$。

（1） 求椭圆 $C$ 的标准方程；

（2） 若动点 $P(x_0,y_0)$ 为椭圆 $C$ 外一点，且点 $P$ 到椭圆 $C$ 的两条切线相互垂直，求点 $P$ 的轨迹方程。

\end{question}
\begin{question}
\gkti{21}已知函数 $f(x)=\dfrac13x^3+x^2+ax+1$（$a\in\mathbf R$）。

（1） 求函数 $f(x)$ 的单调区间；

（2） 当 $a<0$ 时，试讨论是否存在 $x_0\in\left(0,\dfrac12\right)\cup\left(\dfrac12,1\right)$，使得 $f(x_0)=f\left(\dfrac12\right)$。

\end{question}